\section{Reproducibility and artifacts}

The JevBench code is open source under the MIT licence at
\url{https://github.com/Leanmcp/jevbench}. The multilingual extension adds four
scripts to it: \texttt{translate\_cases.py} translates and checks the cases and
records every translated string with its model, reasoning effort and call;
\texttt{45\_run\_multilingual.sh} runs a model on a language;
\texttt{compare\_langs.py} computes the paired comparison; and
\texttt{build\_multilingual.py} produces every number, table and figure in this
paper and on the project website. Translated cases keep the identifiers of their
English originals, and predictions are keyed by case identifier and carry no
source text. The results and all examples are browsable at
\url{https://leanmcp.github.io/jevbench/multilingual.html}.

\section*{Ethics statement}

The benchmark contains harmful, offensive and manipulative text by design,
taken from public safety datasets, and we translated it into eight languages
faithfully, without softening it. Translated harmful text can be misused in the
same ways as the English originals; it adds no harmful capability that a public
translation service does not already provide. Our finding that guards miss
harmful requests in some languages is a vulnerability report. We describe it at
the level of languages and models, and give no recipe beyond what prior work on
multilingual jailbreaks has published.

\section*{Use of AI assistance}

The translations were produced by a language model, as described in
Section~\ref{sec:translation}. The authors also used AI assistance for software
engineering, drafting and language editing of this manuscript. All experimental
design, methodological decisions and interpretation of results are the
authors' own. Every number in the paper is generated from the released runs by
the scripts above, and the authors checked the cited works against their
official records. The authors take full responsibility for the content.
